\chapter{Giới thiệu}

\section{Giới thiệu đề tài}

Nền tảng video ngắn (short-video) như TikTok tạo ra lượng sự kiện tương tác khổng lồ theo thời gian thực: xem, bỏ qua nhanh, thích, chia sẻ, song song với dòng video mới được các nhà sáng tạo nội dung (Provider) đăng tải liên tục. Hệ thống gợi ý phải đồng thời giải quyết ba bài toán: tận dụng lịch sử dài hạn (collaborative filtering), phản ứng gần như tức thời với hành vi ngắn hạn (streaming), và phân phối được nội dung mới (item cold-start) trên một giao diện tiêu thụ mượt mà.

Trong bài tập lớn môn \textbf{Dữ liệu lớn (CO3137)}, nhóm xây dựng \textbf{StreamRec} --- hệ thống mô phỏng nền tảng gợi ý video ngắn theo mô hình \textbf{Providers \& Customers (like TikTok)}, triển khai trên \textbf{kiến trúc Lambda}:

\begin{itemize}
  \item \textbf{Batch layer:} học điểm collaborative filtering dài hạn bằng \textbf{Apache Spark MLlib ALS} chạy định kỳ tự động trên cụm \texttt{spark-master:7077} (kèm đường fallback Truncated SVD nhẹ trong API khi cần).
  \item \textbf{Speed layer:} xử lý dòng sự kiện \texttt{events.user} từ \textbf{Apache Kafka} theo cửa sổ tumbling 60\,giây để cập nhật sở thích ngắn hạn (\path{user:affinity:{user_id}}) và xu hướng (\path{metrics:hot_tags}, \path{metrics:hot_tags:1m}) trên \textbf{Redis}.
  \item \textbf{Cấu trúc dữ liệu xác suất (Probabilistic Data Structures):} sử dụng \textbf{Bloom Filter} trên Redis để lọc video đã xem (seen-filtering) cho từng người dùng, và \textbf{HyperLogLog} (\texttt{PFADD}\slash\texttt{PFCOUNT}) để ước lượng số người dùng hoạt động duy nhất (unique active users) với bộ nhớ cố định $\approx 12$\,KB.
  \item \textbf{Serving layer:} FastAPI hợp nhất điểm $\alpha\cdot\text{content} + \beta\cdot\text{CF} + \gamma\cdot\text{realtime}$, bổ sung \textbf{Cold-start Traffic Pool (Recency Boost)} cho video mới đăng và \textbf{Diversity Re-ranking} tránh lặp liên tiếp cùng creator/chủ đề.
\end{itemize}

\section{Đặc trưng 4V của Big Data trong bài toán}

Hệ thống StreamRec được thiết kế để phản ánh đầy đủ bốn đặc trưng \textbf{4V} của Dữ liệu lớn ở quy mô bài tập lớn:

\begin{center}
{\small\SingleSpacing
\renewcommand{\arraystretch}{1.12}
\begin{tabularx}{\textwidth}{@{}>{\raggedright\arraybackslash}p{2.6cm}>{\raggedright\arraybackslash}X@{}}
\toprule
Đặc trưng & Thể hiện cụ thể trong StreamRec \\
\midrule
\textbf{Volume} & 1\,000 video thật ($\approx 11{,}3$\,GB), tập tương tác gốc 795\,071 dòng (231\,258 dòng sau khi lọc theo các video có file \texttt{.mp4}). \\
\textbf{Velocity} & Dòng sự kiện tương tác được \texttt{traffic-simulator} phát lại theo timeline thực (tăng tốc \texttt{SPEEDUP}, đạt trần tới 200\,events/giây) qua Kafka, kết hợp \texttt{upload-simulator} đẩy 5 video mới mỗi 15\,giây. \\
\textbf{Variety} & Kết hợp dữ liệu phi cấu trúc (video \texttt{.mp4}, ảnh thumbnail \texttt{.jpg}), bán cấu trúc (JSON event trên Kafka, metadata \texttt{ffprobe}), và có cấu trúc (PostgreSQL, CSV 28 cột). \\
\textbf{Veracity} & Xử lý dữ liệu nhiễu/thiếu: lọc noise tag (\texttt{tag\_filters}), kiểm tra tính toàn vẹn file media trước khi xếp hạng (\texttt{media\_file\_ok}), giải quyết cold-start cho cả user mới và video mới. \\
\bottomrule
\end{tabularx}}
\end{center}

\section{Bài toán Providers \& Customers và dataset ShortVideo}

Bài toán gồm hai vai trò tách biệt (được kiểm chứng trực tiếp trên toàn bộ \path{interaction_filtered.csv} --- ví dụ trong 50\,000 dòng đầu có 1\,095 \texttt{user\_id} và 851 \texttt{author\_id} với độ giao nhau bằng 0):

\begin{itemize}
  \item \textbf{Customer (người xem):} sử dụng giao diện TikTok Web gồm \textbf{For~You} (feed dọc toàn màn hình, tự phát hiện quick-skip $<2$\,s, tính tỷ lệ xem \textit{completion rate}) và \textbf{Explore} (khám phá video theo xu hướng Speed Layer).
  \item \textbf{Provider (người đăng):} sử dụng \textbf{TikTok Studio} để đăng video kèm preview trực tiếp; video đi qua pipeline bất đồng bộ \texttt{draft} $\rightarrow$ \texttt{extracting} $\rightarrow$ \texttt{enriched} rồi vào Cold-start Traffic Pool.
\end{itemize}

Nhóm \textbf{không crawl} TikTok / YouTube / Kuaishou production. Nguồn dữ liệu chính là bộ dữ liệu học thuật \textbf{Tsinghua ShortVideo} (giấy phép CC BY 4.0) gồm file \texttt{.mp4} thật, bảng tương tác 28 cột và metadata (category, title). Dữ liệu tương tác sau khi lọc theo 1\,000 video có file thật (231\,258 dòng) được \textbf{tách theo mốc thời gian \texttt{exposed\_time}}: 70\% sớm (161\,886 dòng) nuôi Batch CF; 30\% muộn (69\,372 dòng) để \texttt{traffic-simulator} replay lên Kafka. Tương tự, 1\,000 video được chia thành 700 video catalog ấm và 300 video trong \texttt{upload\_pool} cho \texttt{upload-simulator}.

\section{Mục tiêu}

\begin{itemize}
  \item Xây dựng pipeline thu thập, tách dữ liệu theo thời gian và nạp dữ liệu ShortVideo thật vào stack Docker Compose.
  \item Triển khai kiến trúc Lambda với Kafka làm event bus, Speed Layer xử lý cửa sổ tumbling 60\,s, và Spark Batch ALS chạy định kỳ trên cụm Spark.
  \item Ứng dụng cấu trúc dữ liệu xác suất (Bloom Filter và HyperLogLog) vào đúng bài toán thực tế của nền tảng video ngắn.
  \item Nâng cấp Hybrid Recommendation tiệm cận hành vi TikTok (completion rate, quick-skip, seen-filtering, cold-start boost, diversity re-ranking) và đánh giá định lượng offline (Precision@K, Recall@K, NDCG@K) trên tập chia 70/30.
  \item Xây dựng giao diện TikTok Web thuần túy kèm nút \textbf{Show Big Data} ẩn/hiện theo ngữ cảnh từng trang và trang \textbf{System Ops (/dashboard)} tách riêng.
\end{itemize}
